% Teil II: der durchgerechnete Weg.

\subsection*{Gegeben und gesucht}
Gegeben: Öldichte $\rho=\qty{900}{kg/m^3}$ (aus dem Blatt),
Höhendifferenz $\Delta h=\qty{0.8}{cm}$ auf $L=\qty{10}{m}$,
Tiefe $d=\qty{6}{cm}$, $g=\qty{9.81}{m/s^2}$. Die Viskosität $\eta$ des Öls
ist auf dem eingescannten Blatt \emph{nicht lesbar} (siehe Merksatz unten).
Gesucht: (a) Bedeutung von $\mathrm{Re}$; (b) $\mathrm v(d)$; (c)
$\mathrm{Re}$ und laminar/turbulent.

\subsection*{Ansatz}
(a) Konzeptfrage. (b) Gleichgewicht Schubspannung $\leftrightarrow$
Hangabtrieb für eine Schicht. (c) Mittelwert + Definitionsformel, Zahlen
soweit verfügbar.

\subsection*{Rechnung}
\textbf{a)}\quad$\mathrm{Re}=\rho\bar v d/\eta$ charakterisiert das Verhältnis
Trägheits- zu Viskosenkräften und damit den Übergang laminar/turbulent.

\textbf{b) Profil und Oberflächengeschwindigkeit.} Für eine Schicht der Dicke
$d-y$ mit Grundfläche $A$ gilt im stationären Fall: Reibungskraft
$F_R=\eta A \frac{\mathrm dv}{\mathrm dy}$ gleich der Hangabtriebskraft
$\rho g A(d-y)\sin\alpha$. Kürzen von $A$, auflösen und integrieren
($v(0)=0$):
\begin{equation}
  \eta\frac{\mathrm dv}{\mathrm dy}=\rho g (d-y)\sin\alpha
  \;\Rightarrow\;
  v(y)=\frac{\rho g\sin\alpha}{\eta}\Big(d\,y-\frac{y^2}{2}\Big)\;.
\end{equation}
An der Oberfläche $y=d$:
\begin{equation}
  v(d)=\frac{\rho g\sin\alpha\, d^2}{2\,\eta}\;.
\end{equation}
(macht die geforderte Formel plausibel.) Neigung:
$\sin\alpha\approx\frac{0.008}{10}=8\times10^{-4}$.

\textbf{c) Reynoldszahl.} Mittelwert des Parabelprofils:
\begin{equation}
  \bar v=\frac{2}{3}v(d)
  =\frac{\rho g\sin\alpha\, d^2}{3\,\eta}\;,\qquad
  \mathrm{Re}=\frac{\rho\,\bar v\,d}{\eta}
  =\frac{\rho^2 g\sin\alpha\, d^3}{3\,\eta^2}\;.
\end{equation}
Mit den gegebenen Werten:
\begin{equation}
  \mathrm{Re}
  =\frac{(900)^2\cdot 9.81\cdot 8\times10^{-4}\cdot(0.06)^3}{3\,\eta^2}
  \approx \frac{4.57\times10^{4}\ \mathrm{Pa^2\,s^{-2}}}{\eta^2}\;.
\end{equation}
Ohne den Zahlenwert der Viskosität $\eta$ lässt sich $\mathrm{Re}$ nicht
abschliessend bewerten; für typische Öle ($\eta\approx10^{-1}\,$--$10^1\,
\mathrm{Pa\,s}$) ergäbe sich laminare Strömung — das ist aber die
anzugebende Zahlenaussage.

\begin{ergebnis}
\begin{equation*}
  v(d)=\frac{\rho g\sin\alpha\, d^2}{2\eta}\;,\quad
  \bar v=\frac23 v(d)\;,\quad
  \mathrm{Re}=\frac{\rho\,\bar v\,d}{\eta}
  =\frac{\rho^2 g\sin\alpha\, d^3}{3\eta^2}\;.
\end{equation*}
Der finale Zahlenwert von $\mathrm{Re}$ (und damit die laminar/turbulent-
Aussage) benötigt $\eta$.
\end{ergebnis}

\begin{kontrolle}
\textbf{Einheiten:} $v(d)$:
$\frac{kg/m^3\cdot m/s^2\cdot m^2}{Pa\cdot s}
=\frac{kg\,m^2/s^2}{kg/(m\,s^2)\,\cdot s}=m/s$. \checkmark

\textbf{Querprobe:} $v(d)\propto\sin\alpha$ — stärkere Neigung, schneller;
$\bar v=\frac23 v(d)$ korrekt für Parabel (nicht für lineares Profil).

\textbf{Größenordnung:} Mit $d=6\,\mathrm{cm}$ und kleiner Neigung erwartet
man Oberflächengeschwindigkeiten im Bereich weniger
$\mathrm{cm/s}$ bei typischer Ölviskosität — langsame, stratifizierte
Strömung.
\end{kontrolle}

\begin{merksatz}
\textbf{Nicht lesbarer Wert:} Die Viskosität $\eta$ für Aufgabe A4 ist auf
dem eingescannten Blatt nicht lesbar (OCR konnte sie nicht auflösen). Die
Formeln und der Rechenweg sind vollständig; setze den Zahlenwert von $\eta$
aus dem Original ein, um $\mathrm{Re}$ zu einem konkreten Wert zu bringen.
Die allgemeine Aussage bleibt: kleine Neigung und hohe Viskosität
$\Rightarrow$ laminar.
\end{merksatz}
